\section{Introduction}\label{sec:intro}

Between 2021 and 2026 Britain experienced four trade events in close
sequence: three acts of trade policy (a new customs border with its
largest trading partner, 2021; a tariff war with its largest
single-country export market, 2025; and a new trade agreement with
India, 2026) and a large trade-transmitted price movement, the 2022
terms-of-trade deterioration. The
policy/market distinction is carried explicitly throughout: the energy
episode is not a policy act, and the comparison never treats it as
one. Each episode has its own literature, and several have been studied at
household level: the TCA's food-price effect by income decile, the
energy shock's incidence and relief in detail. What is missing is the
join: none has been put through the statutory tax--benefit system that
determines what households finally keep, and none has been compared
with the others under a common lens. This paper supplies the join in
full for the three price-led episodes and the tariff episode's
earnings channel, and the comparative lens,
with its conditions declared, for all five.

The design holds the second stage fixed and varies the shock. Each
episode enters as a \emph{declared first stage} (published estimates
for the negative shocks, the government's own projections for the
positive ones, never re-estimated here) and is mapped onto households
through two channels: consumer prices against expenditure patterns,
and labour income through the statutory tax--benefit system. Because
gross sizes span two orders of magnitude, comparisons run in ratios:
burden as a share of spending by income decile and cushioning shares. And because the rulebook itself moved
across the window (the Covid uplift removed in 2021; upratings lagging
a year behind prices; the Energy Price Guarantee legislated
mid-shock), the design treats the tax--benefit vintage as an explicit
dimension: the same shock is scored under the rules of its day, under
the previous year's rules, and with the discretionary layer switched
off.

The first-pass results already qualify the received wisdom in three
ways. First, the negative shocks were regressive and the positive
episodes were not progressive: on its realised path the energy shock
took \MsEnergyRealisedPctSpendDecOne\ per cent of bottom-decile
spending against \MsEnergyRealisedPctSpendDecTen\ per cent at the
top. The TCA food-price effect was
\MsTcaWindowAvgPctSpendDecOne\ against
\MsTcaWindowAvgPctSpendDecTen\ per cent on the window-average path,
while the India agreement's consumer gains are proportional and worth
a few pounds a year, too small to register as a share of spending.
Second, each episode's regressivity is generated by its category's
budget-share gradient: food takes \MsTcaBudgetSharePctDecOne\ per
cent of bottom-decile spending against \MsTcaBudgetSharePctDecTen\
at the top, energy \MsEnergyBudgetSharePctDecOne\ against
\MsEnergyBudgetSharePctDecTen, clothing and footwear roughly equal
shares everywhere: composition, not size, decides who bears a trade
shock. Third, the finding the rulebook axis exists for: the modelled
responses run through different instruments on different margins. The
earnings-led tariff scenario was cushioned automatically
(\MsUSTariffDisplacementCushionPct\ to
\MsUSTariffWageCutCushionPct\ per cent, computed in this
repository); the
price-led energy shock was cushioned by discretion, the Guarantee
alone taking \MsEnergyCushionShareFyPct\ per cent of the
counterfactual shock; and indexation, the price side's one automatic
channel, followed the previous September's inflation, leaving a
single adult's Universal Credit allowance
\pounds\MsUcLagCostGbp\ below contemporaneous indexation in the
crisis year, months after the \pounds\MsUcUpliftRemovalGbp\ uplift
ended. The two margins carry different estimands, as the anatomy
table's notes make explicit.

The contribution is threefold: a statutory household study of four
UK trade events, three price-led episodes through the tax and benefit
system and the tariff episode through the earnings channel, with the
conditions every episode's numbers carry stated throughout.
First, the statutory second stage itself. The trade literature
measures shocks without the system that redistributes them; the
microsimulation literature measures the system against income shocks,
for which it is designed. Joining them for the energy and post-Brexit
food episodes shows that, with nominal incomes fixed, the modelled cash benefits do
not respond to a \emph{price} shock (the model's own Energy Price
Guarantee variable does, and the paper counts it as discretionary),
that the discretionary response to one shock was large and cushioned
lower deciles at higher rates while the other received no instrument,
and that, on imputed exposure, most of the variance in the energy
burden lies within income deciles rather than between them
(Section~\ref{sec:results-secondstage}); a closing exercise
aggregates the measured incidence into MPC-weighted first-round
demand and Atkinson-weighted welfare under declared parameters
(Section~\ref{sec:results-mpc}). Second, the rulebook axis: holding a shock fixed and
varying the vintage of the welfare state that receives it. It is executed here for the
largest episode, the same households scored under the 2022 and 2023
parameter systems, which prices the crisis-year uprating lag through
the full statutory structure (Section~\ref{sec:results-rulebook}),
and specified for the rest (Section~\ref{sec:roadmap}); post-Covid
Britain, four non-benchmark episodes landing on four rulebooks, is
its natural laboratory. Third, the comparative set, offered as context rather than as the
delivered analysis: the tariff episode is carried as a declared
earnings scenario, every first stage is labelled estimate or
projection, and every dial is declared. The comparative design cannot rank episodes
on welfare and its cross-episode ratios are descriptive; it records
that, on an episode-specific attribution, the energy shock was met
with energy-specific instruments and the food price rise with none.

The paper proceeds in reading order. Section~\ref{sec:framework}
defines the estimand, states what does and does not require
identification, sets out the four episodes with their declared first
stages, and describes the public-data first pass.
Section~\ref{sec:results} reports the first pass;
Section~\ref{sec:results-secondstage} runs the price-led episodes,
the energy shock, the TCA food price rise and the India agreement,
through the statutory system on enhanced-FRS household records in
PolicyEngine UK, and the tariff episode through the earnings-channel
pipeline, where
the automatic-versus-discretionary decomposition, the
rulebook-vintage run, the within-decile dispersion and the demand and
welfare aggregation are computed. Section~\ref{sec:discussion}
concludes and specifies the extension to the remaining episodes
(Section~\ref{sec:roadmap}).
